\section{Użycie LLM w projekcie}

% TODO: Wymagane przez prowadzącego. Pokryj szczegółowo:
%  1. Lista zadań delegowanych do Claude Opus 4.7 (np. szkielet repo, silnik gry,
%     warianty MCTS, harness eksperymentów, GUI, skrypty analityczne, raport)
%  2. Każdy fragment podlegał review autorów --- udokumentowane w commitach git
%  3. Bugi wykryte w wygenerowanym kodzie i sposób ich naprawy (np. issue z
%     wraparoundem przesunięć, validacja from_string, overflow w u8 multiplication)
%  4. Użycie LLM do raportu: pomoc w strukturze, formułowaniu, ale wszystkie
%     interpretacje wyników i twierdzenia merytoryczne pochodzą od autorów
%  5. Świadomość ograniczeń: weryfikacja statystyczna nadal po stronie autorów

\subsection{Zakres użycia}
[Co dokładnie zostało wygenerowane przez LLM, co napisali autorzy.]

\subsection{Proces weryfikacji}
[Kod review, testy cross-validation, code review LLM]

\subsection{Wykryte błędy}
[Lista bugów wykrytych w wygenerowanym kodzie i jak je naprawiono.]

\subsection{Świadomość ograniczeń}
[Zasady używania LLM, ostatnia decyzja zawsze należy do autorów.]
